\scanpage{566}
\mspage{47\deleted{\unsure{7}}6 bis}
L'idée, c'est qu'on appellerait (sans hypothèse
de petitesse) \uline{catégorie accessible} $\unsure{\mathcal E|E}$ sur une
autre $C$, une catégorie ayant la propriété
qu'on vient de dire --- mais en supposant
\inserted{\unsure{sans doute|toutefois}} de plus que $C$ satisfasse à $L_\pi$, j'imagine :
pour tt foncteur accessible
$M\mathrel{\unsure{-|\longrightarrow}}C$, avec $M$ accessible,
$\mathcal E\times_C M$ est accessible et le foncteur
canonique $\mathcal E\times_C M\longrightarrow\mathcal E$ est accessible.

\deleted{4.17. \uline{Retour sur les \deleted{\unsure{opérations}}\inserted{\unsure{paratopos}} et les
opérations} $M\boxtimes N,\underline{\operatorname{Hom}}_{!}(M,N)$.}

\deleted{4.16.7. Retour sur 4.16.1 :

\uline{Théorème} 4.16.7.1. Soit $M$ une catégorie accessible.
Alors a) Toute sous-catégorie \uline{fermée} $\underline S$ de $M$
est accessible dans $M$, et b) L'ensemble
de ces sous-catégories est \uline{petit}.}

\deleted{\uline{Corollaire} 4.16.7.2. Soit $M$ une catégorie accessible.
Pour que toute sous-catégorie ouverte $\unsure{\underline T|\overline{\underline I}|\overline{\underline J}}$ de
$M$ soit accessible, il f. et s. qu'il en
soit ainsi des \deleted{\unsure{catég}} sous-catégories $M_z$
considérées dans 4.16.1., i.e. (pour $z\in\operatorname{Ob}M$)}
% archive sheet number: 565
